\section{Positive coefficient, realization and no-wrap}
\label{sec:positive-realization}

The Fourier analysis is complete before this section begins.  We now use its positive exact coefficient to obtain a subset in the quotient, translate the congruence into a rational identity modulo one, and only then apply the independent load bound that rules out wrapping.

\begin{theorem}[Avoiding parameterized fixed target]
\label{thm:avoiding-fixed-target}
Let $t=a/b\in(0,1)$ be reduced with squarefree $b$, let $\eta\in(0,1)$ be fixed, let $\gamma$ satisfy \cref{eq:admissible-region}, and let $\mathcal T$ be finite.  For every sufficiently large $X$, there is a subset $A\subseteq\mathcal E_X$, disjoint from $\mathcal T$, such that
\[
 \sum_{e\in A}\frac1e=t.
\]
\end{theorem}

\begin{proof}
Let $W$ be the Bernoulli weight of subsets $A\subseteq\mathcal E$ satisfying
\[
 \sum_{e\in A}\frac{L}{e}\equiv\frac{aL}{b}\pmod L.
\]
By \cref{eq:reciprocal-fourier,thm:positive-fourier-numerator},
\[
 LW=\sum_{h\bmod L}F(h)
\]
has positive real part.  The left side is real and nonnegative, so $W>0$.  Hence some subset satisfies the target congruence.  Dividing by $L$, its reciprocal sum differs from $t$ by an integer.

By \cref{eq:exact-centering-load},
\[
 0\leq\sum_{e\in A}\frac1e\leq\Lambda<1,
 \qquad
 0<t<1.
\]
Thus the integer difference lies strictly between $-1$ and $1$ and must be zero.  The constructed denominators are distinct squarefree semiprimes and avoid $\mathcal T$ by \cref{prop:arithmetic-capacity}.
\end{proof}

The quotient witness, the rational congruence and the no-wrap conclusion are therefore three distinct statements.  The load interval is used only after positivity, and the exact identity is not part of the abstract composition theorem.  With the direct target theorem established, the characterization follows by a separate prime-dilution argument.
